\section{Introduction}
\label{sec:intro}

We ask whether knowing where a shortcut acts inside a language model tells us where to repair it. To test this, we
plant a shortcut with known ground truth in GPT-2 small, locate it with three interpretability analyses, and compare
that location with where retraining a single block, and a training-free read-out, repair it. We find that the repair
data, not the location, decides whether retraining a block passes our mitigation criterion, while the location affects
how fast a block repairs the shortcut and by what mechanism.

Shortcut learning~\cite{geirhos2020shortcut} lets a model reach high accuracy through a cue that does not generalize.
Interpretability offers tools to \emph{find} where such a behaviour acts inside a network: reading intermediate
predictions with the logit lens~\cite{nostalgebraist2020logitlens,belrose2023tunedlens}, swapping components between
models, or patching activations~\cite{vig2020causal,meng2022locating,heimersheim2024patching}. It is natural to assume
that the place these tools point to is also where a \emph{fix} should go. For factual knowledge, Hase et
al.~\cite{hase2023localization} showed that this assumption fails: causal tracing does not predict which layer is best
to edit, and work on unlearning reached related conclusions~\cite{lee2025unlearning,guo2025mechanistic}. It is less
clear what happens for a shortcut picked up during fine-tuning, and why.

Our shortcut is an invented actor name that always comes with the same label in a sentiment classifier's training
data, so we know exactly what the model should and should not use. We locate it with the logit lens, block swaps and
activation patching (Sec.~\ref{sec:find}), retrain single blocks with and without counter-examples in the repair data
(Secs.~\ref{sec:fix} and~\ref{sec:how}), read the answer out of an intermediate layer (Sec.~\ref{sec:readout}), and
test how far the results carry (Sec.~\ref{sec:generalize}).

\textbf{Contributions.} (i)~Converging evidence from block swaps and activation patching, with the logit lens, that
in full fine-tuning the name's information stays at the name tokens until blocks 9--11 move it to the answer.
(ii)~A controlled comparison showing that, for single-block retraining, the counter-examples in the repair data decide
whether a block passes: with them every single block passes, without them none does (under full fine-tuning and LoRA
rank~64; at rank~4 only some blocks pass even with them). (iii)~Evidence that location sets the speed of a partial
repair, with two fast places, block~0 and the later blocks; our analyses single out the later blocks but not
block~0. The pattern was found post hoc and confirmed on new seeds by a registered test; patching shows how block~0
does it. (iv)~A training-free read-out that works best at the
located layer, which we treat as a consistency check because it repeats the lens measurement. (v)~Checks of how far
this carries: LoRA fine-tuning, a second model (Pythia-160M: find analyses, and retraining with counter-examples only), a word-choice cue that the tools place
much earlier, the name on a second dataset (Yelp, retraining only), and memorized wrong labels, which have no shared cue to contradict (Sec.~\ref{sec:generalize}). Every prediction we registered and its outcome is listed in
Table~\ref{tab:prereg} or Table~\ref{tab:prereg-rest}.
